\documentclass{article}
\usepackage[margin=28mm]{geometry}
\usepackage{amsmath,amssymb}
\usepackage{xurl}
\usepackage[hidelinks]{hyperref}
% Shared commands for notes in docs/paper-gaps/.

\DeclareUnicodeCharacter{2080}{\ifmmode{}_{0}\else\textsubscript{0}\fi}
\DeclareUnicodeCharacter{2081}{\ifmmode{}_{1}\else\textsubscript{1}\fi}
\DeclareUnicodeCharacter{2082}{\ifmmode{}_{2}\else\textsubscript{2}\fi}
\DeclareUnicodeCharacter{2099}{\ifmmode{}_{n}\else\textsubscript{n}\fi}

\newcommand{\Lean}{\textsc{Lean}}
\ExplSyntaxOn
\NewDocumentCommand{\leanid}{m}{
  \ifmmode
    \textsf{\detokenize{#1}}
  \else
    \group_begin:
    \urlstyle{sf}
    \tl_set:Nn \l_tmpa_tl {#1}
    \tl_replace_all:Nnn \l_tmpa_tl {\_} {_}
    \exp_args:NV \path \l_tmpa_tl
    \group_end:
  \fi
}
\ExplSyntaxOff
\providecommand{\tr}{\operatorname{tr}}
\newcommand{\tnleanIssueBase}{https://github.com/LionSR/TNLean/issues/}
\newcommand{\tnleanPrBase}{https://github.com/LionSR/TNLean/pull/}
\newcommand{\tnleanDocBase}{https://sirui-lu.com/TNLean/docs/}
\newcommand{\ghissue}[1]{\href{\tnleanIssueBase#1}{GitHub issue \##1}}
\newcommand{\ghpr}[1]{\href{\tnleanPrBase#1}{PR \##1}}
\newcommand{\doclink}[1]{\href{\tnleanDocBase#1.html}{\path{#1}}}

% Dirac notation and the identity operator, as in blueprint/src/macros/common.tex.
% \providecommand keeps note-local definitions authoritative.
\usepackage{dsfont}
\providecommand{\ket}[1]{|#1\rangle}
\providecommand{\bra}[1]{\langle #1|}
\providecommand{\braket}[2]{\langle #1 | #2 \rangle}
\providecommand{\ketbra}[2]{|#1\rangle\!\langle #2|}
\providecommand{\Id}{\mathds{1}}
\providecommand{\one}{\mathds{1}}
\providecommand{\C}{\mathbb{C}}
\providecommand{\MN}[1]{M_{#1}(\C)}
\providecommand{\MD}{M_D(\C)}

% Verdict marker: \gapnote{<kind>}{<status>}. Kind is the policy
% classification (clarification, local-correction, scope-restriction,
% unfaithful, false-source, open-gap); status is open, wip (elimination
% underway), resolved, or historical. Machine-read by the site index;
% prints a verdict line.
\newcommand{\gapnote}[2]{\par\noindent\textbf{Verdict:} #1 (#2).\par}

\providecommand{\gapnote}[2]{\par\noindent\textbf{Verdict:} #1 (#2).\par}
\providecommand{\leanid}[1]{\textsf{\detokenize{#1}}}
\title{The Number of Factors in Hastings' Error Bound at Every Temperature}
\author{}
\date{}
\begin{document}
\maketitle
\gapnote{local-correction}{resolved}

\section{The source statement}
Hastings, \emph{Solving gapped Hamiltonians locally}~\cite{Hastings2006Solving},
\href{https://arxiv.org/abs/cond-mat/0508554v3}{arXiv:cond-mat/0508554v3},
source \path{proj3new.tex}, lines 1303--1314, sets
$\rho(\beta,l_{\mathrm{proj}})=\rho(\beta/n,l_{\mathrm{proj}})^n$ for
$\beta>\beta_0$ (eq.~(mplt)), with $n$ the smallest even integer such that
$\beta/n\le\beta_0$. Lines 1334--1341 bound the operator $\rho_m$ of the terms
with at least $m$ large clusters across the $n$ factors by
\begin{align}
  Z^{-1}\operatorname{Tr}(O\rho_m)\le\frac1{m!}
  \Bigl(\sum_{\mathrm{clusters}}ny^{n_C}\Bigr)^m\lVert O\rVert,
  \label{eq:hastings05_temperature_factor_count_rhom}
\end{align}
and lines 1342--1347 conclude, ``following the same arguments as before'',
\begin{align}
  Z^{-1}\bigl|\operatorname{Tr}[O\rho(\beta)]
  -\operatorname{Tr}[O\rho(\beta,l_{\mathrm{proj}})]\bigr|
  \le\bigl\{\exp[(\beta/\beta_0)Ne^{-l_{\mathrm{proj}}/\xi_{\mathrm{perc}}}]
  -1\bigr\}\lVert O\rVert.
  \label{eq:hastings05_temperature_factor_count_printed}
\end{align}

\section{The point at issue}
Summing (\ref{eq:hastings05_temperature_factor_count_rhom}) over $m\ge1$ and
bounding $\sum_Cy^{n_C}$ by $Ne^{-l_{\mathrm{proj}}/\xi_{\mathrm{perc}}}$ as
at lines 1227--1254 gives
\begin{align}
  \bigl\{\exp[nNe^{-l_{\mathrm{proj}}/\xi_{\mathrm{perc}}}]-1\bigr\}
  \lVert O\rVert,
  \label{eq:hastings05_temperature_factor_count_n}
\end{align}
with the number $n$ of factors in the exponent. The minimality of $n$ gives
$\beta/\beta_0\le n<\beta/\beta_0+2$, so $n$ may exceed $\beta/\beta_0$ by
almost $2$, and the printed bound
(\ref{eq:hastings05_temperature_factor_count_printed}) is smaller than
(\ref{eq:hastings05_temperature_factor_count_n}). The source's argument does
not give the printed bound; the factor $\beta/\beta_0$ stands for $n$.

\section{The statement proved}
The formalization proves (\ref{eq:hastings05_temperature_factor_count_n}) for
every number $n\ge1$ of factors. For the source's $n$ this is at most
$\{\exp[(\beta/\beta_0+2)Ne^{-l_{\mathrm{proj}}/\xi_{\mathrm{perc}}}]-1\}
\lVert O\rVert$. The conclusion the source draws, that $l_{\mathrm{proj}}$
growing logarithmically in the system size suffices for a good approximation
(lines 1348--1349), is unchanged.

\section{Formalization}
The bound (\ref{eq:hastings05_temperature_factor_count_n}) is
\leanid{TNLean.Lattice.ClusterExpansion.norm_trace_mul_exp_sub_truncatedExp_pow_le_percolation},
derived from
\leanid{TNLean.Lattice.ClusterExpansion.norm_trace_mul_exp_sub_truncatedExp_pow_le}.
There the size threshold counts the terms of a large connected set rather
than its diameter, as recorded in the note on cluster size and diameter in
the percolation bound.
\bibliographystyle{unsrt}
\bibliography{references}
\end{document}
